\documentclass[fleqn]{article}
\usepackage{amsmath}
\DeclareMathOperator{\sen}{sen} 
\usepackage{amsfonts}
\usepackage{enumitem}
\title{Solución Problemas Fourier Continuo}
\author{Juan Rodríguez}

\setlength{\parindent}{0pt}

\begin{document}
\maketitle
\section{Ejercicio 1}
	Determinar si las siguientes señales son pares o impares:

	\[
	x_1(t)=t\cos(t)
	\]

	\[
	x_2(t)=\cos(t)\sin^2(t)
	\]

	\[
	x_3(t)=t\sin(t)
	\]

	\subsection{Solución}

	Para $x_1(t)$:

	\[
	x_1(-t)
	=
	(-t)\cos(-t)
    =
	-t\cos(t)
	=
	-x_1(t)
	\]

	Por tanto:

	\[
	\boxed{
	x_1(t)\text{ es impar}
	}
	\]

	Para $x_2(t)$:

	\[
	x_2(-t)
	=
	\cos(-t)\sin^2(-t)
	=
	\cos(t)\sin^2(t)
	=
	x_2(t)
	\]

	Por tanto:

	\[
	\boxed{
	x_2(t)\text{ es par}
	}
	\]

	Para $x_3(t)$:

	\[
	x_3(-t)
	=
	(-t)\sin(-t)
	=
	(-t)(-\sin(t))
	=
	t\sin(t)
	=
	x_3(t)
	\]

	Por tanto:

	\[
	\boxed{
	x_3(t)\text{ es par}
	}
	\]

\section{Ejercicio 2}
	Determinar la parte par y la parte impar de la señal

	\[
	x(t)=2\cos(t)-\sen(t)+3\sen(t)\cos(t)
	\]

	\subsection{Solución}

	La parte par viene dada por:

	\[
	x_{\text{par}}(t)
	=
	\frac{x(t)+x(-t)}{2}
	\]

	Calculamos:

	\[
	x(-t)
	=
	2\cos(-t)-\sen(-t)+3\sen(-t)\cos(-t)
	\]

	\[
	=
	2\cos(t)+\sen(t)-3\sen(t)\cos(t)
	\]

	Por tanto:

	\[
	x_{\text{par}}(t)
	=
	\frac{
	2\cos(t)-\sen(t)+3\sen(t)\cos(t)
	+
	2\cos(t)+\sen(t)-3\sen(t)\cos(t)
	}{2}
	\]

	\[
	\boxed{
	x_{\text{par}}(t)=2\cos(t)
	}
	\]

	La parte impar viene dada por:

	\[
	x_{\text{impar}}(t)
	=
	\frac{x(t)-x(-t)}{2}
	\]

	\[
	x_{\text{impar}}(t)
	=
	\frac{
	2\cos(t)-\sen(t)+3\sen(t)\cos(t)
	-
	\left(
	2\cos(t)+\sen(t)-3\sen(t)\cos(t)
	\right)
	}{2}
	\]

	\[
	\boxed{
	x_{\text{impar}}(t)
	=
	-\sen(t)+3\sen(t)\cos(t)
	}
	\]

\section{Ejercicio 3}
	Obtener el período de las señales

	\[
	x_1(t)=\cos\left(\frac{8\pi}{31}t\right)
	\]

	y

	\[
	x_2(t)=\sen\left(\frac{2}{5}t\right)
	\]

	\subsection{Solución}

	Para una señal sinusoidal:

	\[
	T_0=\frac{2\pi}{\omega_0}
	\]

	Para $x_1(t)$:

	\[
	\omega_0=\frac{8\pi}{31}
	\]

	\[
	T_1
	=
	\frac{2\pi}{\frac{8\pi}{31}}
	=
	\frac{31}{4}
	\]

	Por tanto:

	\[
	\boxed{
	T_1=\frac{31}{4}
	}
	\]

	Para $x_2(t)$:

	\[
	\omega_0=\frac{2}{5}
	\]

	\[
	T_2
	=
	\frac{2\pi}{\frac{2}{5}}
	=
	5\pi
	\]

	Por tanto:

	\[
	\boxed{
	T_2=5\pi
	}
	\]

\section{Ejercicio 4}
	Determinar si la señal

	\[
	x(t)=\cos(3.5\pi t)+\sen(2\pi t)
	+2\cos\left(\frac{7\pi}{6}t\right)
	\]

	es periódica y, en caso afirmativo, determinar su período fundamental.

	\subsection{Solución}

	Calculamos el período de cada término.

	Para:

	\[
	\cos(3.5\pi t)
	=
	\cos\left(\frac{7\pi}{2}t\right)
	\]

	se tiene:

	\[
	T_1
	=
	\frac{2\pi}{\frac{7\pi}{2}}
	=
	\frac{4}{7}
	\]

	Para:

	\[
	\sen(2\pi t)
	\]

	se tiene:

	\[
	T_2
	=
	\frac{2\pi}{2\pi}
	=
	1
	\]

	Para:

	\[
	2\cos\left(\frac{7\pi}{6}t\right)
	\]

	se tiene:

	\[
	T_3
	=
	\frac{2\pi}{\frac{7\pi}{6}}
	=
	\frac{12}{7}
	\]

	Buscamos el menor valor $T_0$ que sea múltiplo de los tres períodos:

	\[
	T_0
	=
	12
	\]

	En efecto:

	\[
	12
	=
	21T_1
	=
	12T_2
	=
	7T_3
	\]

	Por tanto, la señal es periódica y su período fundamental es:

	\[
	\boxed{
	T_0=12
	}
	\]

\section{Ejercicio 5}
	Determinar si son periódicas las señales

	\[
	x_1(t)=\cos(3\pi t)+\sen(6t)
	\]

	y

	\[
	x_2(t)=\cos(5t)+\sen\left(\frac{4}{3}t\right)
	\]

	y, en caso afirmativo, determinar su período fundamental.

	\subsection{Solución}

	Para que la suma de dos señales sinusoidales sea periódica, el cociente entre sus frecuencias angulares debe ser racional.

	Para $x_1(t)$:

	\[
	\omega_1=3\pi,
	\qquad
	\omega_2=6
	\]

	\[
	\frac{\omega_1}{\omega_2}
	=
	\frac{3\pi}{6}
	=
	\frac{\pi}{2}
	\]

	Como:

	\[
	\frac{\pi}{2}\notin\mathbb{Q}
	\]

	la señal no es periódica.

	\[
	\boxed{
	x_1(t)\text{ no es periódica}
	}
	\]

	Para $x_2(t)$:

	\[
	\omega_1=5,
	\qquad
	\omega_2=\frac{4}{3}
	\]

	\[
	\frac{\omega_1}{\omega_2}
	=
	\frac{5}{\frac{4}{3}}
	=
	\frac{15}{4}
	\]

	Como:

	\[
	\frac{15}{4}\in\mathbb{Q}
	\]

	la señal es periódica.

	Calculamos los períodos:

	\[
	T_1
	=
	\frac{2\pi}{5}
	\]

	\[
	T_2
	=
	\frac{2\pi}{\frac{4}{3}}
	=
	\frac{3\pi}{2}
	\]

	Buscamos el menor múltiplo común:

	\[
	15T_1
	=
	4T_2
	=
	6\pi
	\]

	Por tanto:

	\[
	\boxed{
	T_0=6\pi
	}
	\]

\section{Ejercicio 6}
	Calcular las siguientes expresiones:

	\[
	\text{a) }\int_{-\infty}^{\infty}e^{-t}\delta(t)\,dt
	\]

	\[
	\text{b) }\int_{-\infty}^{\infty}e^{-t}\delta(t-3)\,dt
	\]

	\[
	\text{c) }\int_{-\infty}^{\infty}(t^2-2)\delta(t)\,dt
	\]

	\[
	\text{d) }\int_{-\infty}^{\infty}e^{-t}\delta(2t)\,dt
	\]

	\subsection{Solución}

	Utilizamos la propiedad de selectividad:

	\[
	\int_{-\infty}^{\infty}
	f(t)\delta(t-t_0)\,dt
	=
	f(t_0)
	\]

	Para el apartado a):

	\[
	\int_{-\infty}^{\infty}
	e^{-t}\delta(t)\,dt
	=
	e^0
	\]

	\[
	\boxed{1}
	\]

	Para el apartado b):

	\[
	\int_{-\infty}^{\infty}
	e^{-t}\delta(t-3)\,dt
	=
	e^{-3}
	\]

	\[
	\boxed{e^{-3}}
	\]

	Para el apartado c):

	\[
	\int_{-\infty}^{\infty}
	(t^2-2)\delta(t)\,dt
	=
	0^2-2
	\]

	\[
	\boxed{-2}
	\]

	Para el apartado d), utilizamos:

	\[
	\delta(at)
	=
	\frac{1}{|a|}\delta(t)
	\]

	Por tanto:

	\[
	\delta(2t)
	=
	\frac{1}{2}\delta(t)
	\]

	\[
	\int_{-\infty}^{\infty}
	e^{-t}\delta(2t)\,dt
	=
	\frac{1}{2}
	\int_{-\infty}^{\infty}
	e^{-t}\delta(t)\,dt
	\]

	\[
	\boxed{\frac{1}{2}}
	\]

\section{Ejercicio 7}
	Calcular la energía y la potencia media de la señal

	\[
	x(t)=2
	\]

	\subsection{Solución}

	La energía de una señal es:

	\[
	E
	=
	\int_{-\infty}^{\infty}
	|x(t)|^2\,dt
	\]

	Como:

	\[
	|x(t)|^2=4
	\]

	se tiene:

	\[
	E
	=
	\int_{-\infty}^{\infty}4\,dt
	=
	+\infty
	\]

	Por tanto:

	\[
	\boxed{
	E=+\infty
	}
	\]

	La potencia media es:

	\[
	P
	=
	\lim_{T\to\infty}
	\frac{1}{2T}
	\int_{-T}^{T}
	|x(t)|^2\,dt
	\]

	\[
	P
	=
	\lim_{T\to\infty}
	\frac{1}{2T}
	\int_{-T}^{T}
	4\,dt
	\]

	\[
	P
	=
	\lim_{T\to\infty}
	\frac{1}{2T}
	(8T)
	\]

	\[
	\boxed{
	P=4
	}
	\]

	Por tanto, la señal está definida en potencia:

	\[
	\boxed{
	E=+\infty,
	\qquad
	P=4
	}
	\]

\section{Ejercicio 8}
	Calcular la energía de la señal

	\[
	x(t)=e^{-2t}u(t)
	\]

	\subsection{Solución}

	Como:

	\[
	u(t)=
	\begin{cases}
	0, & t<0 \\[1mm]
	1, & t>0
	\end{cases}
	\]

	la señal queda:

	\[
	x(t)=
	\begin{cases}
	0, & t<0 \\[1mm]
	e^{-2t}, & t>0
	\end{cases}
	\]

	La energía es:

	\[
	E
	=
	\int_{-\infty}^{\infty}
	|x(t)|^2\,dt
	\]

	\[
	E
	=
	\int_0^{\infty}
	e^{-4t}\,dt
	\]

	\[
	E
	=
	\left[
	-\frac{1}{4}e^{-4t}
	\right]_0^{\infty}
	\]

	\[
	E
	=
	0-\left(-\frac{1}{4}\right)
	\]

	Por tanto:

	\[
	\boxed{
	E=\frac{1}{4}
	}
	\]

\section{Ejercicio 9}
	Calcular la energía y la potencia media de la señal

	\[
	x(t)=e^{-|t|}
	\]

	\subsection{Solución}

	La energía es:

	\[
	E
	=
	\int_{-\infty}^{\infty}
	|x(t)|^2\,dt
	\]

	\[
	E
	=
	\int_{-\infty}^{\infty}
	e^{-2|t|}\,dt
	\]

	Como la función es par:

	\[
	E
	=
	2\int_0^{\infty}
	e^{-2t}\,dt
	\]

	\[
	E
	=
	2
	\left[
	-\frac{1}{2}e^{-2t}
	\right]_0^{\infty}
	\]

	\[
	E=1
	\]

	Por tanto:

	\[
	\boxed{
	E=1
	}
	\]

	La potencia media es:

	\[
	P
	=
	\lim_{T\to\infty}
	\frac{1}{2T}
	\int_{-T}^{T}
	e^{-2|t|}\,dt
	\]

	Como la función es par:

	\[
	P
	=
	\lim_{T\to\infty}
	\frac{1}{T}
	\int_0^T
	e^{-2t}\,dt
	\]

	\[
	P
	=
	\lim_{T\to\infty}
	\frac{1}{T}
	\left[
	-\frac{1}{2}e^{-2t}
	\right]_0^T
	\]

	\[
	P
	=
	\lim_{T\to\infty}
	\frac{1-e^{-2T}}{2T}
	=
	0
	\]

	Por tanto:

	\[
	\boxed{
	P=0
	}
	\]

	La señal está definida en energía:

	\[
	\boxed{
	E=1,
	\qquad
	P=0
	}
	\]

\section{Ejercicio 10}
	Calcular la potencia media de la señal

	\[
	x(t)=A\cos(2\pi f_0t+\phi)
	+\frac{A}{2}\cos(6\pi f_0t)
	\]

	\subsection{Solución}

	La señal es periódica. Sus frecuencias son:

	\[
	f_1=f_0,
	\qquad
	f_2=3f_0
	\]

	Por tanto, el período fundamental es:

	\[
	T_0=\frac{1}{f_0}
	\]

	La potencia media es:

	\[
	P
	=
	\frac{1}{T_0}
	\int_0^{T_0}
	|x(t)|^2\,dt
	\]

	Desarrollamos:

	\[
	|x(t)|^2
	=
	A^2\cos^2(2\pi f_0t+\phi)
	+
	\frac{A^2}{4}\cos^2(6\pi f_0t)
	+
	A^2\cos(2\pi f_0t+\phi)\cos(6\pi f_0t)
	\]

	Al integrar en un período, el producto de cosenos de distintas frecuencias tiene integral nula:

	\[
	\int_0^{T_0}
	\cos(2\pi f_0t+\phi)\cos(6\pi f_0t)\,dt
	=
	0
	\]

	Además:

	\[
	\frac{1}{T_0}
	\int_0^{T_0}
	\cos^2(2\pi f_0t+\phi)\,dt
	=
	\frac{1}{2}
	\]

	y:

	\[
	\frac{1}{T_0}
	\int_0^{T_0}
	\cos^2(6\pi f_0t)\,dt
	=
	\frac{1}{2}
	\]

	Por tanto:

	\[
	P
	=
	\frac{A^2}{2}
	+
	\frac{A^2}{4}\frac{1}{2}
	\]

	\[
	P
	=
	\frac{A^2}{2}
	+
	\frac{A^2}{8}
	\]

	\[
	\boxed{
	P=\frac{5A^2}{8}
	}
	\]

\section{Ejercicio 11}
	Dibujar las gráficas de las señales

	\[
	x_1(t)=u(t+3),
	\qquad
	x_2(t)=u(t-2),
	\qquad
	x_3(t)=x_1(t)-x_2(t)
	\]

	\subsection{Solución}

	La señal $x_1(t)=u(t+3)$ es un escalón desplazado $3$ unidades hacia la izquierda:

	\[
	x_1(t)=
	\begin{cases}
	0, & t<-3 \\[1mm]
	\frac{1}{2}, & t=-3 \\[1mm]
	1, & t>-3
	\end{cases}
	\]

	La señal $x_2(t)=u(t-2)$ es un escalón desplazado $2$ unidades hacia la derecha:

	\[
	x_2(t)=
	\begin{cases}
	0, & t<2 \\[1mm]
	\frac{1}{2}, & t=2 \\[1mm]
	1, & t>2
	\end{cases}
	\]

	Finalmente:

	\[
	x_3(t)
	=
	u(t+3)-u(t-2)
	\]

	Por tanto:

	\[
	\boxed{
	x_3(t)=
	\begin{cases}
	0, & t<-3 \\[1mm]
	\frac{1}{2}, & t=-3 \\[1mm]
	1, & -3<t<2 \\[1mm]
	\frac{1}{2}, & t=2 \\[1mm]
	0, & t>2
	\end{cases}
	}
	\]

	Así, $x_3(t)$ es un pulso rectangular de altura $1$ entre $t=-3$ y $t=2$.

\section{Ejercicio 12}
	Dibujar la gráfica de la señal

	\[
	x(t)=3u(t)+tu(t)-(t-1)u(t-1)-5u(t-2)
	\]

	\subsection{Solución}

	Estudiamos la señal por intervalos.

	Si $t<0$:

	\[
	u(t)=u(t-1)=u(t-2)=0
	\]

	\[
	x(t)=0
	\]

	Si $0<t<1$:

	\[
	u(t)=1,
	\qquad
	u(t-1)=u(t-2)=0
	\]

	\[
	x(t)=3+t
	\]

	Si $1<t<2$:

	\[
	u(t)=u(t-1)=1,
	\qquad
	u(t-2)=0
	\]

	\[
	x(t)
	=
	3+t-(t-1)
	\]

	\[
	x(t)=4
	\]

	Si $t>2$:

	\[
	u(t)=u(t-1)=u(t-2)=1
	\]

	\[
	x(t)
	=
	3+t-(t-1)-5
	\]

	\[
	x(t)=-1
	\]

	Por tanto:

	\[
	\boxed{
	x(t)=
	\begin{cases}
	0, & t<0 \\[1mm]
	3+t, & 0<t<1 \\[1mm]
	4, & 1<t<2 \\[1mm]
	-1, & t>2
	\end{cases}
	}
	\]

	La gráfica parte de $0$, salta a $3$ en $t=0$, crece linealmente hasta $4$ en $t=1$, permanece constante en $4$ hasta $t=2$ y después baja a $-1$.

\section{Ejercicio 13}
	Dada la señal $x(t)$ de la figura, obtener las gráficas de

	\[
	x(t-2),\quad x(t+1),\quad x(2t),\quad x(0.1t),
	\quad x(-t),\quad x(-t-2),\quad x(-t+3)
	\]

	\subsection{Solución}

	La señal original puede escribirse como:

	\[
	x(t)=
	\begin{cases}
	-1, & -1\leq t<0 \\[1mm]
	\displaystyle\frac{t}{2}, & 0\leq t\leq2 \\[1mm]
	0, & \text{resto}
	\end{cases}
	\]

	Para $x(t-2)$, desplazamos la señal $2$ unidades hacia la derecha:

	\[
	\boxed{
	x(t-2)=
	\begin{cases}
	-1, & 1\leq t<2 \\[1mm]
	\displaystyle\frac{t-2}{2}, & 2\leq t\leq4 \\[1mm]
	0, & \text{resto}
	\end{cases}
	}
	\]

	Para $x(t+1)$, desplazamos la señal $1$ unidad hacia la izquierda:

	\[
	\boxed{
	x(t+1)=
	\begin{cases}
	-1, & -2\leq t<-1 \\[1mm]
	\displaystyle\frac{t+1}{2}, & -1\leq t\leq1 \\[1mm]
	0, & \text{resto}
	\end{cases}
	}
	\]

	Para $x(2t)$, comprimimos la señal por un factor $2$:

	\[
	\boxed{
	x(2t)=
	\begin{cases}
	-1, & -\frac{1}{2}\leq t<0 \\[1mm]
	t, & 0\leq t\leq1 \\[1mm]
	0, & \text{resto}
	\end{cases}
	}
	\]

	Para $x(0.1t)$, expandimos la señal por un factor $10$:

	\[
	\boxed{
	x(0.1t)=
	\begin{cases}
	-1, & -10\leq t<0 \\[1mm]
	\displaystyle\frac{t}{20}, & 0\leq t\leq20 \\[1mm]
	0, & \text{resto}
	\end{cases}
	}
	\]

	Para $x(-t)$, invertimos la señal respecto del eje vertical:

	\[
	\boxed{
	x(-t)=
	\begin{cases}
	\displaystyle-\frac{t}{2}, & -2\leq t\leq0 \\[1mm]
	-1, & 0<t\leq1 \\[1mm]
	0, & \text{resto}
	\end{cases}
	}
	\]

	Para $x(-t-2)$:

	\[
	\boxed{
	x(-t-2)=
	\begin{cases}
	\displaystyle\frac{-t-2}{2}, & -4\leq t\leq-2 \\[1mm]
	-1, & -2<t\leq-1 \\[1mm]
	0, & \text{resto}
	\end{cases}
	}
	\]

	Para $x(-t+3)$:

	\[
	\boxed{
	x(-t+3)=
	\begin{cases}
	\displaystyle\frac{3-t}{2}, & 1\leq t\leq3 \\[1mm]
	-1, & 3<t\leq4 \\[1mm]
	0, & \text{resto}
	\end{cases}
	}
	\]

\section{Ejercicio 14}
	Obtener el desarrollo en serie de Fourier de la señal

	\[
	x(t)=\cos(2\pi t)-\sen(3\pi t)
	\]

	\subsection{Solución}

	Las frecuencias angulares son:

	\[
	\omega_1=2\pi,
	\qquad
	\omega_2=3\pi
	\]

	Buscamos una frecuencia fundamental $\omega_0$ tal que:

	\[
	2\pi=2\omega_0
	\]

	y:

	\[
	3\pi=3\omega_0
	\]

	Por tanto:

	\[
	\omega_0=\pi
	\]

	y el período fundamental es:

	\[
	T_0=\frac{2\pi}{\omega_0}=2
	\]

	La serie de Fourier tiene la forma:

	\[
	x(t)
	=
	a_0
	+
	\sum_{k=1}^{\infty}
	a_k\cos(k\omega_0t)
	+
	\sum_{k=1}^{\infty}
	b_k\sen(k\omega_0t)
	\]

	Comparando con:

	\[
	x(t)
	=
	\cos(2\pi t)-\sen(3\pi t)
	\]

	se obtiene:

	\[
	a_0=0
	\]

	\[
	a_2=1
	\]

	\[
	b_3=-1
	\]

	y el resto de coeficientes son nulos:

	\[
	a_k=0
	\qquad
	k\neq2
	\]

	\[
	b_k=0
	\qquad
	k\neq3
	\]

	Por tanto:

	\[
	\boxed{
	\omega_0=\pi,
	\qquad
	T_0=2
	}
	\]

	y el desarrollo en serie de Fourier es:

	\[
	\boxed{
	x(t)
	=
	\cos(2\omega_0t)
	-
	\sen(3\omega_0t)
	}
	\]

	con:

	\[
	\boxed{
	a_0=0,\qquad
	a_2=1,\qquad
	b_3=-1
	}
	\]

\section{Ejercicio 15}
	Obtener el desarrollo en serie de Fourier de la señal

	\[
	x(t)
	=
	1+\frac{1}{2}\cos(2\pi t)
	+\cos(4\pi t)
	+\frac{2}{3}\cos(6\pi t)
	\]

	\subsection{Solución}

	Las frecuencias angulares son:

	\[
	\omega_1=2\pi,
	\qquad
	\omega_2=4\pi,
	\qquad
	\omega_3=6\pi
	\]

	Por tanto, la frecuencia fundamental es:

	\[
	\omega_0=2\pi
	\]

	y el período fundamental:

	\[
	T_0
	=
	\frac{2\pi}{\omega_0}
	=
	1
	\]

	La serie de Fourier tiene la forma:

	\[
	x(t)
	=
	a_0
	+
	\sum_{k=1}^{\infty}
	a_k\cos(k\omega_0t)
	+
	\sum_{k=1}^{\infty}
	b_k\sen(k\omega_0t)
	\]

	Comparando términos:

	\[
	a_0=1
	\]

	\[
	a_1=\frac{1}{2}
	\]

	\[
	a_2=1
	\]

	\[
	a_3=\frac{2}{3}
	\]

	y:

	\[
	a_k=0,
	\qquad
	k\geq4
	\]

	Como no aparecen términos seno:

	\[
	b_k=0
	\qquad
	\forall k
	\]

	Por tanto:

	\[
	\boxed{
	\omega_0=2\pi,
	\qquad
	T_0=1
	}
	\]

	y los coeficientes son:

	\[
	\boxed{
	a_0=1,\qquad
	a_1=\frac{1}{2},\qquad
	a_2=1,\qquad
	a_3=\frac{2}{3}
	}
	\]

	\[
	\boxed{
	b_k=0
	\qquad
	\forall k
	}
	\]

\section{Ejercicio 16}
	Obtener el desarrollo en serie de Fourier de la señal

	\[
	x(t)
	=
	2+\cos\left(\frac{2\pi}{3}t\right)
	+4\sen\left(\frac{5\pi}{3}t\right)
	\]

	\subsection{Solución}

	Las frecuencias angulares son:

	\[
	\omega_1=\frac{2\pi}{3},
	\qquad
	\omega_2=\frac{5\pi}{3}
	\]

	Buscamos una frecuencia fundamental $\omega_0$ tal que:

	\[
	\frac{2\pi}{3}=2\omega_0
	\]

	y:

	\[
	\frac{5\pi}{3}=5\omega_0
	\]

	Por tanto:

	\[
	\omega_0=\frac{\pi}{3}
	\]

	y el período fundamental es:

	\[
	T_0
	=
	\frac{2\pi}{\omega_0}
	=
	6
	\]

	La serie de Fourier tiene la forma:

	\[
	x(t)
	=
	a_0
	+
	\sum_{k=1}^{\infty}
	a_k\cos(k\omega_0t)
	+
	\sum_{k=1}^{\infty}
	b_k\sen(k\omega_0t)
	\]

	Comparando términos:

	\[
	a_0=2
	\]

	\[
	a_2=1
	\]

	\[
	b_5=4
	\]

	y el resto de coeficientes son nulos:

	\[
	a_k=0
	\qquad
	k\neq2
	\]

	\[
	b_k=0
	\qquad
	k\neq5
	\]

	Por tanto:

	\[
	\boxed{
	\omega_0=\frac{\pi}{3},
	\qquad
	T_0=6
	}
	\]

	y los coeficientes no nulos son:

	\[
	\boxed{
	a_0=2,\qquad
	a_2=1,\qquad
	b_5=4
	}
	\]

\section{Ejercicio 17}
	Obtener el desarrollo en serie de Fourier de la señal

	\[
	x(t)=\sum_{n\in\mathbb{Z}}z(t+4n)
	\]

	donde:

	\[
	z(t)=e^{-t},
	\qquad
	t\in[0,4]
	\]

	\subsection{Solución}

	El período fundamental es:

	\[
	T_0=4
	\]

	Por tanto:

	\[
	\omega_0
	=
	\frac{2\pi}{T_0}
	=
	\frac{\pi}{2}
	\]

	La serie de Fourier tiene la forma:

	\[
	x(t)
	=
	a_0
	+
	\sum_{k=1}^{\infty}
	a_k\cos(k\omega_0t)
	+
	\sum_{k=1}^{\infty}
	b_k\sen(k\omega_0t)
	\]

	Calculamos $a_0$:

	\[
	a_0
	=
	\frac{1}{4}
	\int_0^4e^{-t}\,dt
	\]

	\[
	a_0
	=
	\frac{1}{4}
	\left[
	-e^{-t}
	\right]_0^4
	=
	\frac{1-e^{-4}}{4}
	\]

	Para $a_k$:

	\[
	a_k
	=
	\frac{1}{2}
	\int_0^4
	e^{-t}
	\cos\left(\frac{k\pi}{2}t\right)\,dt
	\]

	\[
	\int
	e^{-t}\cos(\alpha t)\,dt
	=
	\frac{e^{-t}}{1+\alpha^2}
	\left(
	-\cos(\alpha t)+\alpha\sen(\alpha t)
	\right)
	\]

	con:

	\[
	\alpha=\frac{k\pi}{2}
	\]

	Como:

	\[
	\cos(2k\pi)=1,
	\qquad
	\sen(2k\pi)=0
	\]

	se obtiene:

	\[
	a_k
	=
	\frac{1-e^{-4}}
	{2\left(1+\frac{k^2\pi^2}{4}\right)}
	\]

	\[
	\boxed{
	a_k
	=
	\frac{2(1-e^{-4})}
	{4+k^2\pi^2}
	}
	\]

	Para $b_k$:

	\[
	b_k
	=
	\frac{1}{2}
	\int_0^4
	e^{-t}
	\sen\left(\frac{k\pi}{2}t\right)\,dt
	\]

	\[
	\int
	e^{-t}\sen(\alpha t)\,dt
	=
	-\frac{e^{-t}}{1+\alpha^2}
	\left(
	\sen(\alpha t)+\alpha\cos(\alpha t)
	\right)
	\]

	Por tanto:

	\[
	b_k
	=
	\frac{
	\frac{k\pi}{2}(1-e^{-4})
	}{
	2\left(1+\frac{k^2\pi^2}{4}\right)
	}
	\]

	\[
	\boxed{
	b_k
	=
	\frac{k\pi(1-e^{-4})}
	{4+k^2\pi^2}
	}
	\]

	Finalmente:

	\[
	\boxed{
	x(t)
	=
	\frac{1-e^{-4}}{4}
	+
	\sum_{k=1}^{\infty}
	\frac{2(1-e^{-4})}{4+k^2\pi^2}
	\cos\left(\frac{k\pi}{2}t\right)
	+
	\sum_{k=1}^{\infty}
	\frac{k\pi(1-e^{-4})}{4+k^2\pi^2}
	\sen\left(\frac{k\pi}{2}t\right)
	}
	\]

\section{Ejercicio 18}
	Obtener el desarrollo en serie de Fourier de la señal

	\[
	x(t)=\sum_{n\in\mathbb{Z}}z(t+6n)
	\]

	donde:

	\[
	z(t)=\operatorname{sig}(t),
	\qquad
	t\in[-3,3]
	\]

	\subsection{Solución}

	El período fundamental es:

	\[
	T_0=6
	\]

	Por tanto:

	\[
	\omega_0
	=
	\frac{2\pi}{T_0}
	=
	\frac{\pi}{3}
	\]

	Como la señal es impar:

	\[
	a_0=0
	\]

	y:

	\[
	a_k=0
	\]

	Calculamos los coeficientes $b_k$:

	\[
	b_k
	=
	\frac{2}{6}
	\int_{-3}^{3}
	\operatorname{sig}(t)
	\sen\left(\frac{k\pi}{3}t\right)\,dt
	\]

	Como el producto es una función par:

	\[
	b_k
	=
	\frac{2}{3}
	\int_0^3
	\sen\left(\frac{k\pi}{3}t\right)\,dt
	\]

	\[
	b_k
	=
	\frac{2}{3}
	\left[
	-\frac{3}{k\pi}
	\cos\left(\frac{k\pi}{3}t\right)
	\right]_0^3
	\]

	\[
	b_k
	=
	\frac{2}{k\pi}
	\left(
	1-\cos(k\pi)
	\right)
	\]

	Como:

	\[
	\cos(k\pi)=(-1)^k
	\]

	se obtiene:

	\[
	b_k
	=
	\frac{2}{k\pi}
	\left(
	1-(-1)^k
	\right)
	\]

	Por tanto:

	\[
	b_k=
	\begin{cases}
	\displaystyle\frac{4}{k\pi}, & k\text{ impar} \\[2mm]
	0, & k\text{ par}
	\end{cases}
	\]

	\[
	\boxed{
	x(t)
	=
	\frac{4}{\pi}
	\sum_{n=0}^{\infty}
	\frac{1}{2n+1}
	\sen\left(
	\frac{(2n+1)\pi}{3}t
	\right)
	}
	\]

	\[
	\boxed{
	x(t)
	=
	\frac{4}{\pi}
	\left[
	\sen\left(\frac{\pi t}{3}\right)
	+
	\frac{1}{3}\sen(\pi t)
	+
	\frac{1}{5}\sen\left(\frac{5\pi t}{3}\right)
	+\cdots
	\right]
	}
	\]

\section{Ejercicio 19}
	Obtener el desarrollo en serie de Fourier de la señal

	\[
	x(t)=\sum_{n\in\mathbb{Z}}z(t+2n)
	\]

	donde:

	\[
	z(t)=t^2,
	\qquad
	t\in[-1,1]
	\]

	\subsection{Solución}

	El período fundamental es:

	\[
	T_0=2
	\]

	Por tanto:

	\[
	\omega_0
	=
	\frac{2\pi}{T_0}
	=
	\pi
	\]

	Como la señal es par:

	\[
	b_k=0
	\]

	Calculamos $a_0$:

	\[
	a_0
	=
	\frac{1}{2}
	\int_{-1}^{1}
	t^2\,dt
	\]

	Como $t^2$ es par:

	\[
	a_0
	=
	\int_0^1 t^2\,dt
	\]

	\[
	a_0
	=
	\left[
	\frac{t^3}{3}
	\right]_0^1
	=
	\frac{1}{3}
	\]

	Calculamos ahora $a_k$:

	\[
	a_k
	=
	\frac{2}{2}
	\int_{-1}^{1}
	t^2\cos(k\pi t)\,dt
	\]

	Como el integrando es par:

	\[
	a_k
	=
	2
	\int_0^1
	t^2\cos(k\pi t)\,dt
	\]

	Integramos por partes:

	\[
	\int
	t^2\cos(k\pi t)\,dt
	=
	\frac{t^2\sen(k\pi t)}{k\pi}
	-
	\frac{2}{k\pi}
	\int
	t\sen(k\pi t)\,dt
	\]

	De nuevo por partes:

	\[
	\int
	t\sen(k\pi t)\,dt
	=
	-\frac{t\cos(k\pi t)}{k\pi}
	+
	\frac{\sen(k\pi t)}{k^2\pi^2}
	\]

	Por tanto:

	\[
	\int_0^1
	t^2\cos(k\pi t)\,dt
	=
	\frac{2\cos(k\pi)}{k^2\pi^2}
	\]

	Como:

	\[
	\cos(k\pi)=(-1)^k
	\]

	se obtiene:

	\[
	a_k
	=
	\frac{4(-1)^k}{k^2\pi^2}
	\]

	Por tanto:

	\[
	\boxed{
	a_0=\frac{1}{3},
	\qquad
	a_k=\frac{4(-1)^k}{k^2\pi^2},
	\qquad
	b_k=0
	}
	\]

	\[
	\boxed{
	x(t)
	=
	\frac{1}{3}
	+
	\frac{4}{\pi^2}
	\sum_{k=1}^{\infty}
	\frac{(-1)^k}{k^2}
	\cos(k\pi t)
	}
	\]

\section{Ejercicio 20}
	Obtener la transformada de Fourier de las siguientes señales utilizando la definición o alguna propiedad.

	\subsection{Solución}

	\textbf{a) } $x_1(t)=\delta(t)$

	Usando la pareja básica:

	\[
	\boxed{
	\delta(t)
	\overset{\mathcal{F}}{\longleftrightarrow}
	1
	}
	\]

	\textbf{b) } $x_2(t)=1$

	Usando la pareja básica:

	\[
	\boxed{
	1
	\overset{\mathcal{F}}{\longleftrightarrow}
	2\pi\delta(\omega)
	}
	\]

	\textbf{c) } $x_3(t)=e^{j\omega_0t}$

	Usando la pareja básica:

	\[
	\boxed{
	e^{j\omega_0t}
	\overset{\mathcal{F}}{\longleftrightarrow}
	2\pi\delta(\omega-\omega_0)
	}
	\]

	\textbf{d) } $x_4(t)=\cos(\omega_0t)$

	Como:

	\[
	\cos(\omega_0t)
	=
	\frac{
	e^{j\omega_0t}
	+
	e^{-j\omega_0t}
	}{2}
	\]

	por linealidad:

	\[
	X_4(\omega)
	=
	\frac{1}{2}
	\left[
	2\pi\delta(\omega-\omega_0)
	+
	2\pi\delta(\omega+\omega_0)
	\right]
	\]

	Por tanto:

	\[
	\boxed{
	X_4(\omega)
	=
	\pi
	\left[
	\delta(\omega-\omega_0)
	+
	\delta(\omega+\omega_0)
	\right]
	}
	\]

	\textbf{e) }

	\[
	x_5(t)
	=
	\sum_{k=-\infty}^{\infty}
	c_k e^{jk\omega_0t}
	\]

	Aplicando linealidad:

	\[
	X_5(\omega)
	=
	\sum_{k=-\infty}^{\infty}
	c_k
	\mathcal{F}
	\left[
	e^{jk\omega_0t}
	\right]
	\]

	\[
	X_5(\omega)
	=
	2\pi
	\sum_{k=-\infty}^{\infty}
	c_k
	\delta(\omega-k\omega_0)
	\]

	Por tanto:

	\[
	\boxed{
	X_5(\omega)
	=
	2\pi
	\sum_{k=-\infty}^{\infty}
	c_k
	\delta(\omega-k\omega_0)
	}
	\]

	\textbf{f) }

	\[
	x_6(t)
	=
	\sum_{k=-\infty}^{\infty}
	\delta(t-kT_0)
	\]

	Usando la pareja de transformadas del tren de impulsos:

	\[
	\boxed{
	X_6(\omega)
	=
	\frac{2\pi}{T_0}
	\sum_{k=-\infty}^{\infty}
	\delta
	\left(
	\omega-\frac{2\pi k}{T_0}
	\right)
	}
	\]

\section{Ejercicio 21}
	Calcular la transformada de Fourier de la señal

	\[
	x(t)=
	\cos\left(\frac{3\pi}{4}t\right)
	-
	\sen\left(\frac{\pi}{2}t\right)
	\]

	\subsection{Solución}

	Utilizamos las parejas:

	\[
	\cos(\omega_0 t)
	\overset{\mathcal{F}}{\longleftrightarrow}
	\pi
	\left[
	\delta(\omega-\omega_0)
	+
	\delta(\omega+\omega_0)
	\right]
	\]

	y:

	\[
	\sen(\omega_0 t)
	\overset{\mathcal{F}}{\longleftrightarrow}
	\frac{\pi}{j}
	\left[
	\delta(\omega-\omega_0)
	-
	\delta(\omega+\omega_0)
	\right]
	\]

	Para el primer término:

	\[
	\cos\left(\frac{3\pi}{4}t\right)
	\overset{\mathcal{F}}{\longleftrightarrow}
	\pi
	\left[
	\delta\left(\omega-\frac{3\pi}{4}\right)
	+
	\delta\left(\omega+\frac{3\pi}{4}\right)
	\right]
	\]

	Para el segundo término:

	\[
	-\sen\left(\frac{\pi}{2}t\right)
	\overset{\mathcal{F}}{\longleftrightarrow}
	-\frac{\pi}{j}
	\left[
	\delta\left(\omega-\frac{\pi}{2}\right)
	-
	\delta\left(\omega+\frac{\pi}{2}\right)
	\right]
	\]

	Por linealidad:

	\[
	\boxed{
	X(\omega)
	=
	\pi
	\left[
	\delta\left(\omega-\frac{3\pi}{4}\right)
	+
	\delta\left(\omega+\frac{3\pi}{4}\right)
	\right]
	-
	\frac{\pi}{j}
	\left[
	\delta\left(\omega-\frac{\pi}{2}\right)
	-
	\delta\left(\omega+\frac{\pi}{2}\right)
	\right]
	}
	\]

\section{Ejercicio 22}
	Determinar la transformada de Fourier de la señal

	\[
	x(t)=
	\begin{cases}
	\cos(\omega_0 t), & t\in\left[-\frac{T_0}{4},\frac{T_0}{4}\right] \\[2mm]
	0, & \text{resto}
	\end{cases}
	\]

	\subsection{Solución}

	Aplicamos la definición de transformada de Fourier:

	\[
	X(\omega)
	=
	\int_{-\infty}^{\infty}
	x(t)e^{-j\omega t}\,dt
	\]

	Como la señal solo es distinta de cero en
	$\left[-\frac{T_0}{4},\frac{T_0}{4}\right]$:

	\[
	X(\omega)
	=
	\int_{-T_0/4}^{T_0/4}
	\cos(\omega_0 t)e^{-j\omega t}\,dt
	\]

	Utilizamos:

	\[
	\cos(\omega_0 t)
	=
	\frac{
	e^{j\omega_0 t}
	+
	e^{-j\omega_0 t}
	}{2}
	\]

	Entonces:

	\[
	X(\omega)
	=
	\frac{1}{2}
	\int_{-T_0/4}^{T_0/4}
	e^{-j(\omega-\omega_0)t}\,dt
	+
	\frac{1}{2}
	\int_{-T_0/4}^{T_0/4}
	e^{-j(\omega+\omega_0)t}\,dt
	\]

	Usamos:

	\[
	\int_{-a}^{a}
	e^{-j\alpha t}\,dt
	=
	\frac{2\sen(\alpha a)}{\alpha}
	\]

	con:

	\[
	a=\frac{T_0}{4}
	\]

	Por tanto:

	\[
	X(\omega)
	=
	\frac{
	\sen\left(
	(\omega-\omega_0)\frac{T_0}{4}
	\right)
	}{
	\omega-\omega_0
	}
	+
	\frac{
	\sen\left(
	(\omega+\omega_0)\frac{T_0}{4}
	\right)
	}{
	\omega+\omega_0
	}
	\]

	\[
	\boxed{
	X(\omega)
	=
	\frac{
	\sen\left(
	\frac{T_0}{4}(\omega-\omega_0)
	\right)
	}{
	\omega-\omega_0
	}
	+
	\frac{
	\sen\left(
	\frac{T_0}{4}(\omega+\omega_0)
	\right)
	}{
	\omega+\omega_0
	}
	}
	\]

\section{Ejercicio 23}
	Obtener la transformada de Fourier de la señal

	\[
	x(t)=be^{-at}u(t),
	\qquad a>0,\; b>0
	\]

	\subsection{Solución}

	Aplicamos la definición de transformada de Fourier:

	\[
	X(\omega)
	=
	\int_{-\infty}^{\infty}
	x(t)e^{-j\omega t}\,dt
	\]

	Como $u(t)=0$ para $t<0$:

	\[
	X(\omega)
	=
	\int_0^{\infty}
	be^{-at}e^{-j\omega t}\,dt
	\]

	\[
	X(\omega)
	=
	b
	\int_0^{\infty}
	e^{-(a+j\omega)t}\,dt
	\]

	\[
	X(\omega)
	=
	b
	\left[
	-\frac{1}{a+j\omega}
	e^{-(a+j\omega)t}
	\right]_0^{\infty}
	\]

	Como $a>0$:

	\[
	\lim_{t\to\infty}
	e^{-(a+j\omega)t}
	=
	0
	\]

	Por tanto:

	\[
	X(\omega)
	=
	\frac{b}{a+j\omega}
	\]

	\[
	\boxed{
	X(\omega)
	=
	\frac{b}{a+j\omega}
	}
	\]

\section{Ejercicio 24}
	Determinar la transformada de Fourier de la señal

	\[
	x(t)=e^{-|t|}
	\]

	\subsection{Solución}

	Aplicamos la definición:

	\[
	X(\omega)
	=
	\int_{-\infty}^{\infty}
	e^{-|t|}e^{-j\omega t}\,dt
	\]

	Como:

	\[
	e^{-|t|}
	=
	\begin{cases}
	e^t, & t<0 \\[1mm]
	e^{-t}, & t\geq0
	\end{cases}
	\]

	se tiene:

	\[
	X(\omega)
	=
	\int_{-\infty}^{0}
	e^te^{-j\omega t}\,dt
	+
	\int_{0}^{\infty}
	e^{-t}e^{-j\omega t}\,dt
	\]

	\[
	X(\omega)
	=
	\int_{-\infty}^{0}
	e^{(1-j\omega)t}\,dt
	+
	\int_{0}^{\infty}
	e^{-(1+j\omega)t}\,dt
	\]

	Calculamos la primera integral:

	\[
	\int_{-\infty}^{0}
	e^{(1-j\omega)t}\,dt
	=
	\left[
	\frac{e^{(1-j\omega)t}}{1-j\omega}
	\right]_{-\infty}^{0}
	\]

	\[
	=
	\frac{1}{1-j\omega}
	\]

	La segunda integral:

	\[
	\int_{0}^{\infty}
	e^{-(1+j\omega)t}\,dt
	=
	\left[
	-\frac{e^{-(1+j\omega)t}}{1+j\omega}
	\right]_{0}^{\infty}
	\]

	\[
	=
	\frac{1}{1+j\omega}
	\]

	Por tanto:

	\[
	X(\omega)
	=
	\frac{1}{1-j\omega}
	+
	\frac{1}{1+j\omega}
	\]

	\[
	X(\omega)
	=
	\frac{(1+j\omega)+(1-j\omega)}
	{(1-j\omega)(1+j\omega)}
	\]

	\[
	X(\omega)
	=
	\frac{2}{1+\omega^2}
	\]

	\[
	\boxed{
	X(\omega)
	=
	\frac{2}{1+\omega^2}
	}
	\]

\section{Ejercicio 25}
	Obtener la convolución de las señales

	\[
	x(t)=u(t)-u(t-T)
	\]

	e

	\[
	y(t)=be^{-at}u(t),
	\qquad a>0,\; b>0
	\]

	\subsection{Solución}

	La convolución viene dada por:

	\[
	z(t)
	=
	x(t)*y(t)
	=
	\int_{-\infty}^{\infty}
	x(\tau)y(t-\tau)\,d\tau
	\]

	Como:

	\[
	x(\tau)=
	\begin{cases}
	1, & 0\leq\tau\leq T \\[1mm]
	0, & \text{resto}
	\end{cases}
	\]

	y:

	\[
	y(t-\tau)
	=
	be^{-a(t-\tau)}u(t-\tau)
	\]

	se tiene:

	\[
	z(t)
	=
	\int_0^T
	be^{-a(t-\tau)}u(t-\tau)\,d\tau
	\]

	Estudiamos por intervalos.

	Si $t<0$:

	\[
	u(t-\tau)=0
	\]

	para todo $\tau\in[0,T]$, luego:

	\[
	z(t)=0
	\]

	Si $0\leq t<T$:

	\[
	0\leq\tau\leq t
	\]

	\[
	z(t)
	=
	\int_0^t
	be^{-a(t-\tau)}\,d\tau
	\]

	\[
	z(t)
	=
	be^{-at}
	\int_0^t e^{a\tau}\,d\tau
	\]

	\[
	z(t)
	=
	be^{-at}
	\left[
	\frac{e^{a\tau}}{a}
	\right]_0^t
	\]

	\[
	z(t)
	=
	\frac{b}{a}
	\left(
	1-e^{-at}
	\right)
	\]

	Si $t\geq T$:

	\[
	z(t)
	=
	\int_0^T
	be^{-a(t-\tau)}\,d\tau
	\]

	\[
	z(t)
	=
	be^{-at}
	\int_0^T e^{a\tau}\,d\tau
	\]

	\[
	z(t)
	=
	\frac{b}{a}
	e^{-at}
	\left(
	e^{aT}-1
	\right)
	\]

	\[
	z(t)
	=
	\frac{b}{a}
	\left(
	e^{-a(t-T)}-e^{-at}
	\right)
	\]

	Por tanto:

	\[
	\boxed{
	z(t)=
	\begin{cases}
	0, & t<0 \\[2mm]
	\displaystyle
	\frac{b}{a}
	\left(
	1-e^{-at}
	\right),
	& 0\leq t<T \\[3mm]
	\displaystyle
	\frac{b}{a}
	\left(
	e^{-a(t-T)}-e^{-at}
	\right),
	& t\geq T
	\end{cases}
	}
	\]

\section{Ejercicio 26}
	Obtener la convolución de las señales

	\[
	x(t)=-u(t+1)+2u(t)-u(t-1)
	\]

	e

	\[
	y(t)=u(t+2)-u(t-2)
	\]

	\subsection{Solución}

	Las señales pueden escribirse como:

	\[
	x(t)=
	\begin{cases}
	-1, & -1<t<0 \\[1mm]
	1, & 0<t<1 \\[1mm]
	0, & \text{resto}
	\end{cases}
	\]

	y:

	\[
	y(t)=
	\begin{cases}
	1, & -2<t<2 \\[1mm]
	0, & \text{resto}
	\end{cases}
	\]

	La convolución viene dada por:

	\[
	z(t)
	=
	x(t)*y(t)
	=
	\int_{-\infty}^{\infty}
	x(\tau)y(t-\tau)\,d\tau
	\]

	Como $y(t-\tau)=1$ cuando:

	\[
	t-2<\tau<t+2
	\]

	la integral corresponde al área de $x(\tau)$ contenida en dicho intervalo.

	Si $t<-3$:

	\[
	z(t)=0
	\]

	Si $-3\leq t<-2$:

	\[
	z(t)
	=
	-\int_{-1}^{t+2}d\tau
	\]

	\[
	z(t)=-(t+3)
	\]

	Si $-2\leq t<-1$:

	\[
	z(t)
	=
	-\int_{-1}^{0}d\tau
	+
	\int_0^{t+2}d\tau
	\]

	\[
	z(t)
	=
	-1+t+2
	=
	t+1
	\]

	Si $-1\leq t<1$, el intervalo contiene completamente ambas partes de $x(\tau)$:

	\[
	z(t)
	=
	-\int_{-1}^{0}d\tau
	+
	\int_0^1d\tau
	\]

	\[
	z(t)=0
	\]

	Si $1\leq t<2$:

	\[
	z(t)
	=
	\int_{t-2}^{1}x(\tau)\,d\tau
	\]

	\[
	z(t)
	=
	-\int_{t-2}^{0}d\tau
	+
	\int_0^1d\tau
	\]

	\[
	z(t)=t-1
	\]

	Si $2\leq t<3$:

	\[
	z(t)
	=
	\int_{t-2}^{1}d\tau
	\]

	\[
	z(t)=3-t
	\]

	Si $t\geq3$:

	\[
	z(t)=0
	\]

	Por tanto:

	\[
	\boxed{
	z(t)=
	\begin{cases}
	0, & t<-3 \\[1mm]
	-(t+3), & -3\leq t<-2 \\[1mm]
	t+1, & -2\leq t<-1 \\[1mm]
	0, & -1\leq t<1 \\[1mm]
	t-1, & 1\leq t<2 \\[1mm]
	3-t, & 2\leq t<3 \\[1mm]
	0, & t\geq3
	\end{cases}
	}
	\]

\section{Ejercicio 27}
	Obtener la convolución de la señal

	\[
	x(t)
	=
	u\left(t+\frac{T}{2}\right)
	-
	u\left(t-\frac{T}{2}\right)
	\]

	con ella misma, donde $T>0$.

	\subsection{Solución}

	La señal puede escribirse como:

	\[
	x(t)=
	\begin{cases}
	1, & -\frac{T}{2}<t<\frac{T}{2} \\[1mm]
	0, & \text{resto}
	\end{cases}
	\]

	La convolución es:

	\[
	z(t)
	=
	x(t)*x(t)
	=
	\int_{-\infty}^{\infty}
	x(\tau)x(t-\tau)\,d\tau
	\]

	La primera señal es distinta de cero cuando:

	\[
	-\frac{T}{2}<\tau<\frac{T}{2}
	\]

	y la segunda cuando:

	\[
	-\frac{T}{2}<t-\tau<\frac{T}{2}
	\]

	es decir:

	\[
	t-\frac{T}{2}<\tau<t+\frac{T}{2}
	\]

	La convolución es igual a la longitud de solapamiento de ambos intervalos.

	Si $t<-T$:

	\[
	z(t)=0
	\]

	Si $-T\leq t<0$:

	\[
	z(t)
	=
	\int_{-T/2}^{t+T/2}d\tau
	\]

	\[
	z(t)
	=
	t+T
	\]

	Si $0\leq t<T$:

	\[
	z(t)
	=
	\int_{t-T/2}^{T/2}d\tau
	\]

	\[
	z(t)
	=
	T-t
	\]

	Si $t\geq T$:

	\[
	z(t)=0
	\]

	Por tanto:

	\[
	\boxed{
	z(t)=
	\begin{cases}
	0, & t<-T \\[1mm]
	t+T, & -T\leq t<0 \\[1mm]
	T-t, & 0\leq t<T \\[1mm]
	0, & t\geq T
	\end{cases}
	}
	\]

	Equivalentemente:

	\[
	\boxed{
	z(t)=
	\begin{cases}
	T-|t|, & |t|<T \\[1mm]
	0, & |t|\geq T
	\end{cases}
	}
	\]


\end{document}